Use the skymap ``Map-OM01'' to answer the questions below.

% FIGURE NEEDED: full-sky map "Map-OM01" handed out with the 2025 sky-map
% round (2025_SKYMAP paper)

\begin{description}
\item[(OM01.1)] Mark all the objects listed below with a box ($\square$)
  around the object. Label each of your marked objects with the corresponding
  Object No.

  \begin{center}
  \begin{tabular}{cccc}
  Object No. & Object Name & Object No. & Object Name \\
  \hline
  1 & $\beta$ Aur & 5 & $\delta$ Gem \\
  2 & $\delta$ Cep & 6 & $\beta$ CVn \\
  3 & $\delta$ Cnc & 7 & $\alpha$ Lyn \\
  4 & $\delta$ Cet & 8 & $\beta$ Per \\
  \end{tabular}
  \end{center}

\item[(OM01.2)] Mark the positions of the following 6 galaxies from the
  Messier Catalogue using a plus sign ($+$) and label them with their
  corresponding Messier number.

  M32, M51, M74, M81, M94, M101

\item[(OM01.3)] Draw the Ecliptic on the map and label it as ``E''.

\item[(OM01.4)] A total solar eclipse occurred on 1 August 2008. At a certain
  place on Earth the totality occurred at local noon.
  \begin{description}
  \item[(OM01.4a)] Mark the position of the Sun at the time of the eclipse
    with a cross ($\times$) and label it as ``S''.
  \item[(OM01.4b)] Draw the Moon at the appropriate position on the map as
    seen from the same location on 28 July 2008 at local noon, and label it as
    ``M''. The drawing should be of appropriate shape and orientation, but
    need not be to scale. The bright side of the Moon should be shaded.
  \end{description}
\end{description}
